\section{Automatic Error Analysis}

\begin{definition}[Floating-Point Expression Syntax]
Fix $N\in\N$ and $\mathbb F=\mathbb F_{E,M}$. Expressions are
formed according to
\[
  e::=c\mid x_i\mid -e\mid e+e\mid e\times e,
  \qquad c\in\mathbb F,\quad 1\le i\le N.
\]
An expression $e$ is a syntactic object, not a function.
\end{definition}

\begin{definition}[Floating-Point Semantics]
For $e$, let $D_e\subseteq\mathbb F^N$ be the set of input
valuations at which all rounded operations in $e$ yield finite values.
The floating-point semantics is the function
$\llbracket e\rrbracket\colon D_e\to\mathbb F$ defined recursively,
for $\sigma=(\sigma_1,\dots,\sigma_N)\in D_e$, by
\[
\begin{aligned}
\llbracket c\rrbracket(\sigma)&=c,
&\llbracket x_i\rrbracket(\sigma)&=\sigma_i,\\
\llbracket-e\rrbracket(\sigma)&=-\llbracket e\rrbracket(\sigma),
&\llbracket e_1+e_2\rrbracket(\sigma)
  &=\llbracket e_1\rrbracket(\sigma)\oplus\llbracket e_2\rrbracket(\sigma),\\
\llbracket e_1\times e_2\rrbracket(\sigma)
  &=\llbracket e_1\rrbracket(\sigma)\otimes\llbracket e_2\rrbracket(\sigma).
\end{aligned}
\]
For constants and variables, $D_e=\mathbb F^N$; for the other
constructors, the domains are restricted to inputs on which their
subexpressions and the indicated rounded operations are finite.
\end{definition}

\begin{definition}[Real-Valued Semantics]
The real-valued interpretation
$\llbracket e\rrbracket_{\mathrm r}\colon\mathbb F^N\to\R$
is obtained from the same recursive rules by replacing
$\oplus$ and $\otimes$ with exact $+$ and $\times$.
Thus $\llbracket e\rrbracket_{\mathrm r}(\sigma)$ evaluates $e$ over
$\R$ at the same input $\sigma$.
\end{definition}

\begin{example}[Different Semantics]
Let $e=(1+x_1)+(-x_1)$, and assume $2^{M+1}\in\mathbb F$.
At $\sigma_1=2^{M+1}$, ties-to-even yields
\[
  \llbracket e\rrbracket(\sigma)=0,
  \qquad\llbracket e\rrbracket_{\mathrm r}(\sigma)=1.
\]
\end{example}

\begin{definition}[Sound Range and Error Analysis]
Fix the input set
\[
  \Sigma\coloneqq\prod_{i=1}^N([l_i,r_i]\cap\mathbb F)
  \subseteq D_e,\qquad l_i,r_i\in\mathbb F,\quad l_i\le r_i.
\]
This condition excludes overflow at every intermediate operation
for inputs in $\Sigma$. When computing bounds recursively, we also
assume the endpoint rounding operations used by the analysis remain finite.
A sound \emph{range analysis} assigns an interval
$\operatorname{range}(e)$, and a sound \emph{absolute-error analysis}
assigns $\operatorname{err}(e)\ge0$, such that
\[
\begin{aligned}
 \llbracket e\rrbracket(\sigma)&\in\operatorname{range}(e),\\
 |\llbracket e\rrbracket(\sigma)-\llbracket e\rrbracket_{\mathrm r}(\sigma)|
 &\le\operatorname{err}(e)
\end{aligned}
\qquad(\forall\sigma\in\Sigma).
\]
These bounds are computed recursively from the syntax and the
rounding estimates; the fixed input set $\Sigma$ is suppressed
in the notation.
\end{definition}

\begin{definition}[Recursive Range Bounds]
Set
\[
  \operatorname{range}(c)=[c,c],\qquad
  \operatorname{range}(x_i)=[l_i,r_i].
\]
If $\operatorname{range}(e)=[a,b]$, put
$\operatorname{range}(-e)=[-b,-a]$. If
$\operatorname{range}(e_i)=[a_i,b_i]$ for $i=1,2$, put
\[
 \operatorname{range}(e_1+e_2)
 =[\rnd(a_1+a_2),\rnd(b_1+b_2)].
\]
\end{definition}

The sum rule follows from monotonicity: if
$\llbracket e_i\rrbracket(\sigma)\in[a_i,b_i]$, then
\[
 a_1+a_2\le\llbracket e_1\rrbracket(\sigma)
 +\llbracket e_2\rrbracket(\sigma)\le b_1+b_2,
\]
and rounding the endpoints preserves this enclosure.

\begin{definition}[Recursive Absolute-Error Bounds]
Set
\[
 \operatorname{err}(c)=\operatorname{err}(x_i)=0,
 \qquad\operatorname{err}(-e)=\operatorname{err}(e).
\]
If $\operatorname{range}(e_i)=[a_i,b_i]$ for $i=1,2$, define
\[
 \operatorname{err}(e_1+e_2)
 =\operatorname{err}(e_1)+\operatorname{err}(e_2)
 +\tfrac12\ulp(\max\{|a_1+a_2|,|b_1+b_2|\}).
\]
\end{definition}

For the addition rule, put $u_i=\llbracket e_i\rrbracket(\sigma)$
and $v_i=\llbracket e_i\rrbracket_{\mathrm r}(\sigma)$. Then
\[
\begin{aligned}
 |\rnd(u_1+u_2)-(v_1+v_2)|
 &\le |\rnd(u_1+u_2)-(u_1+u_2)|+|u_1-v_1|+|u_2-v_2|\\
 &\le\tfrac12\ulp(u_1+u_2)+\operatorname{err}(e_1)+\operatorname{err}(e_2).
\end{aligned}
\]
Since $u_i\in[a_i,b_i]$ and $\ulp$ is nondecreasing in $|x|$,
the displayed recursion bounds the error.


\begin{definition}[Multiplication Bounds]
Suppose $\operatorname{range}(e_i)=[a_i,b_i]$ for $i=1,2$. Let
\[
  p_-\coloneqq\min\{a_1a_2,a_1b_2,b_1a_2,b_1b_2\},\qquad
  p_+\coloneqq\max\{a_1a_2,a_1b_2,b_1a_2,b_1b_2\}.
\]
Define
\[
  \operatorname{range}(e_1\times e_2)
  =[\rnd(p_-),\rnd(p_+)].
\]
For $i=1,2$, write $U_i=\max\{|a_i|,|b_i|\}$ and
$\eta_i=\operatorname{err}(e_i)$. Set
\[
\begin{aligned}
  \operatorname{err}(e_1\times e_2)
  ={}& U_1\eta_2+U_2\eta_1+\eta_1\eta_2\\
     &+\tfrac12\ulp(\max\{|p_-|,|p_+|\}).
\end{aligned}
\]
\end{definition}

Indeed, for $u_i=\llbracket e_i\rrbracket(\sigma)$ and
$v_i=\llbracket e_i\rrbracket_{\mathrm r}(\sigma)$, the exact
product $u_1u_2$ lies in $[p_-,p_+]$, so monotonicity justifies
the range rule. The error follows from
\[
\begin{aligned}
 |\rnd(u_1u_2)-v_1v_2|
 &\le |\rnd(u_1u_2)-u_1u_2|+|u_1u_2-v_1v_2|\\
 &\le\tfrac12\ulp(\max\{|p_-|,|p_+|\})
 +U_1\eta_2+(U_2+\eta_2)\eta_1.
\end{aligned}
\]

\begin{definition}[Square-Root Extension]
The expression grammar may be extended by $\sqrt e$, on input
valuations for which both interpretations of $e$ are nonnegative:
\[
\begin{aligned}
 \llbracket\sqrt e\rrbracket(\sigma)
 &=\rnd\!\bigl(\sqrt{\llbracket e\rrbracket(\sigma)}\bigr),\\
 \llbracket\sqrt e\rrbracket_{\mathrm r}(\sigma)
 &=\sqrt{\llbracket e\rrbracket_{\mathrm r}(\sigma)}.
\end{aligned}
\]
If $\operatorname{range}(e)=[a,b]\subseteq[0,F_{\max}]$, set
\[
\begin{aligned}
 \operatorname{range}(\sqrt e)
 &=[\rnd(\sqrt a),\rnd(\sqrt b)],\\
 \operatorname{err}(\sqrt e)
 &=\sqrt{\operatorname{err}(e)}+\tfrac12\ulp(\sqrt b).
\end{aligned}
\]
\end{definition}
The range bound follows from monotonicity of square root and
rounding. The error rule follows from, for $u,v\ge0$,
\[
 |\sqrt{u}-\sqrt{v}|\le\sqrt{|u-v|},
\]
followed by the triangle inequality and the half-ULP rounding bound.

\begin{remark}[Limitations of Interval Analysis]
Independent bounds lose dependencies between repeated occurrences
of the same variable. For $e=x_1-x_1$ (i.e., $x_1+(-x_1)$) and
$x_1\in[-1,1]$, interval propagation gives $[-2,2]$ although both
semantics evaluate to $0$. Even a sound error bound can be positive
when the actual rounding error vanishes; soundness does not imply
sharpness.
\end{remark}
